%!TEX program = xelatex
\documentclass[12pt, a4paper]{article}

% ── Engine: XeLaTeX ──────────────────────────────────────────────────────────
\usepackage{fontspec}
\usepackage{polyglossia}
\setdefaultlanguage{russian}
\setotherlanguage{english}

\setmainfont{FreeSerif}
\setsansfont{FreeSans}
\setmonofont{FreeMono}

% ── Math ─────────────────────────────────────────────────────────────────────
\usepackage{amsmath, amssymb, amsthm}
\usepackage{mathtools}
\usepackage{unicode-math}
\setmathfont{Latin Modern Math}

% ── Layout ───────────────────────────────────────────────────────────────────
\usepackage[left=3cm, right=2cm, top=2.5cm, bottom=2.5cm]{geometry}
\usepackage{setspace}
\onehalfspacing
\usepackage{indentfirst}
\setlength{\parindent}{1.25cm}

% ── Tables ───────────────────────────────────────────────────────────────────
\usepackage{booktabs}
\usepackage{array}
\usepackage{tabularx}
\usepackage{multirow}

% ── Misc ─────────────────────────────────────────────────────────────────────
\usepackage{graphicx}
\usepackage{float}
\usepackage{caption}
\usepackage{enumitem}
\usepackage{etoolbox}
\usepackage{hyperref}
\hypersetup{colorlinks=true, linkcolor=black, citecolor=black, urlcolor=blue}
\pretocmd{\section}{\clearpage}{}{}

% ── Theorem environments ─────────────────────────────────────────────────────
\newtheorem{definition}{Определение}
\newtheorem{remark}{Замечание}

% ── Custom commands ───────────────────────────────────────────────────────────
\newcommand{\R}{\mathbb{R}}
\newcommand{\E}{\mathbb{E}}
\newcommand{\given}{\,|\,}
\DeclareMathOperator*{\argmax}{arg\,max}
\DeclareMathOperator*{\clip}{clip}

\setcounter{secnumdepth}{3}
\setcounter{tocdepth}{3}

% ─────────────────────────────────────────────────────────────────────────────
\begin{document}

\begin{titlepage}
	\centering
	
	{\large САНКТ-ПЕТЕРБУРГСКИЙ ГОСУДАРСТВЕННЫЙ УНИВЕРСИТЕТ\par}
	\vspace{0.5cm}
	{\large Искусственный интеллект и наука о данных\par}
	
	\vspace{4cm}
	
	{\large Сергиенко Андрей\par}
	
	\vspace{2cm}
	
	{\Large\bfseries Разработка и исследование агента глубокого обучения с подкреплением для решения задачи дискретного размещения геометрических объектов\par}
	
	\vspace{0.5cm}
	
	{\large Отчёт о прохождении учебной практики\par}
	
	\vfill
	
	\begin{flushright}
		Научный руководитель:\\
		Старший преподаватель кафедры системного программирования\\
		Юрий Александрович Андреев
	\end{flushright}
	
	\vfill
	
	{\large Санкт-Петербург\par}
	{\large 2026\par}
	
\end{titlepage}

% ─────────────────────────────────────────────
\tableofcontents
\newpage
% ─────────────────────────────────────────────

% ════════════════════════════════════════════════════════════════
\section{Введение}
% ════════════════════════════════════════════════════════════════

Обучение с подкреплением (Reinforcement Learning, RL) является одним из наиболее активно развивающихся разделов машинного обучения. Агент взаимодействует со средой, получает обратную связь в виде скалярного сигнала вознаграждения и постепенно вырабатывает стратегию, максимизирующую суммарную дисконтированную награду. Отличительная черта RL - отсутствие учителя с размеченными примерами: агент вынужден самостоятельно исследовать пространство состояний и действий.

Данная работа посвящена разработке и исследованию RL-агента для игры \textit{Block Blast!} - комбинаторной головоломки на поле $8 \times 8$, в которой игрок последовательно размещает случайные геометрические фигуры с целью очищать заполненные строки и столбцы. Задача обладает рядом свойств, делающих её нетривиальной для алгоритмов RL:

\begin{itemize}
  \item \textbf{Большое пространство действий}: до $192$ валидных ходов на каждом шаге при ненулевой доле невалидных действий, которые необходимо явно исключать.
  \item \textbf{Разреженное вознаграждение}: очистка линий - редкое событие, большинство шагов не несут немедленного обучающего сигнала.
  \item \textbf{Необходимость планирования внутри раунда}: каждый раунд агент получает три фигуры и должен размещать их последовательно, принимая во внимание доступность позиций для оставшихся фигур.
  \item \textbf{Случайность среды}: следующая тройка фигур стохастична, что исключает \\ детерминированную оптимальную стратегию.
\end{itemize}

Целью работы является выявление условий и степени превосходства обученного агента над эвристическими подходами, установление обстоятельств, при которых такое превосходство достигается, а также исследование влияния информационного содержания пространства наблюдений на качество обучения агента. Для достижения поставленной цели была реализована эволюция архитектуры агента через четыре поколения с последовательным наращиванием сложности тензора состояния. Четвёртое поколение масштабирует задачу на поле $16\times16$ и вводит двухфазный пайплайн обучения: предварительное supervised-обучение на демонстрациях эвристики (Behavior Cloning) с последующим дообучением методом PPO. Параллельно проведено сравнение архитектур: многослойного перцептрона (MLP), свёрточной нейронной сети (CNN) и трансформера (Vision Transformer, ViT~\cite{dosovitskiy2020}).

В качестве бейзлайна использовались случайный агент и эвристический алгоритм, вдохновлённый работой Dellacherie по игре Tetris~\cite{dellacherie2003}. Превзойти эвристику по среднему и стабильности является главным критерием успеха.

% ════════════════════════════════════════════════════════════════
\section{Постановка задачи}
% ════════════════════════════════════════════════════════════════

\subsection{Описание среды}

Block Blast разворачивается на квадратном поле $\mathcal{B} \in \{0, 1\}^{8 \times 8}$. В начале каждого \textit{раунда} агенту выдаётся тройка фигур из пула, включающего все мономино, домино, тримино, тетромино и L-образные фигуры с учётом всех поворотов. Агент должен разместить все три фигуры, выбирая каждый раз одну из $n_{\text{valid}}$ допустимых позиций. Если после размещения строка или столбец оказываются полностью заполненными, они очищаются, и агент получает вознаграждение. Эпизод завершается, когда хотя бы одну из оставшихся фигур невозможно разместить.

\subsection{Марковский процесс принятия решений}

Задача формализована как \textbf{Марковский процесс принятия решений} (Markov Decision Process, MDP)~\cite{sutton2018}:
\begin{equation}
  \mathcal{M} = \langle \mathcal{S},\; \mathcal{A},\; \mathcal{T},\; R,\; \gamma \rangle,
\end{equation}
где:
\begin{itemize}
  \item $\mathcal{S}$ - пространство состояний (поле $8\times 8$ и текущие фигуры);
  \item $\mathcal{A}$ - пространство действий (пара $\langle$фигура, позиция$\rangle$);
  \item $\mathcal{T}: \mathcal{S} \times \mathcal{A} \to \Delta(\mathcal{S})$ - вероятностная функция переходов;
  \item $R: \mathcal{S} \times \mathcal{A} \to \R$ - функция вознаграждения;
  \item $\gamma = 0.99$ - коэффициент дисконтирования.
\end{itemize}

Агент стремится найти политику $\pi_\theta: \mathcal{S} \to \Delta(\mathcal{A})$ (параметризованную вектором $\theta$), максимизирующую суммарную дисконтированную награду:
\begin{equation}
  J(\theta) = \mathbb{E}_{\pi_\theta}\!\left[G_t\right],
  \qquad
  G_t = \sum_{k=0}^{\infty} \gamma^k\, r_{t+k}.
\end{equation}

\subsection{Пространство действий и маска}

Пространство действий имеет размерность $|\mathcal{A}| = 3 \times 8 \times 8 = 192$: каждое из трёх гнёзд для фигур может быть размещено в любую из $64$ позиций поля (позиция задаётся верхним левым углом). Значительная часть действий в каждом состоянии невалидна (фигура выходит за границы или перекрывается с занятыми клетками).

Маска допустимых действий $m: \mathcal{A} \times \mathcal{S} \to \{0, 1\}$ используется для исключения невалидных ходов из распределения политики:
\begin{equation}
  \tilde{\pi}_\theta(a \given s)
  = \frac{\pi_\theta(a \given s)\cdot m(a,s)}{\displaystyle\sum_{a' \in \mathcal{A}} \pi_\theta(a' \given s)\cdot m(a',s)}.
  \label{eq:masking}
\end{equation}

Подход~\eqref{eq:masking} принципиально отличается от штрафа за невалидные действия: градиент от заведомо невозможных ходов равен нулю, что устраняет информационный шум и ускоряет сходимость~\cite{huang2020}.

% ════════════════════════════════════════════════════════════════
\section{Алгоритм обучения}
% ════════════════════════════════════════════════════════════════

\subsection{Proximal Policy Optimization (PPO)}

В качестве алгоритма обучения выбран \textbf{PPO}~\cite{schulman2017} - метод семейства ``policy gradient'', позволяющий делать несколько обновлений политики по одному набору собранных данных, ограничивая при этом величину отклонения от предыдущей политики.

Функция потерь PPO:
\begin{equation}
  L^{\text{CLIP}}(\theta)
  = \mathbb{E}_t\!\left[
      \min\!\left(
        \rho_t(\theta)\,\hat{A}_t,\;
        \clip\!\left(\rho_t(\theta),\,1-\varepsilon,\,1+\varepsilon\right)\hat{A}_t
      \right)
    \right],
  \label{eq:ppo}
\end{equation}
где
\begin{equation}
  \rho_t(\theta) = \frac{\pi_\theta(a_t \given s_t)}{\pi_{\theta_\text{old}}(a_t \given s_t)}, \qquad \varepsilon = 0.2.
\end{equation}

Здесь $\hat{A}_t$ - оценка функции преимущества (advantage), вычисляемая методом \textbf{GAE}~\cite{schulman2015}:
\begin{equation}
  \hat{A}_t^{\text{GAE}(\lambda)}
  = \sum_{l=0}^{\infty} (\gamma\lambda)^l \delta_{t+l},
  \qquad
  \delta_t = r_t + \gamma\, V(s_{t+1}) - V(s_t),
\end{equation}
где $\lambda = 0.95$ - параметр сглаживания, $V(s)$ - функция ценности, оцениваемая отдельной головой сети (критиком).

Полная целевая функция включает потери критика и энтропийный бонус:
\begin{equation}
  L(\theta) = L^{\text{CLIP}}(\theta)
            - c_v\, L^{\text{VF}}(\theta)
            + c_e\, H\!\left[\pi_\theta(\cdot \given s)\right],
\end{equation}
где $c_v = 1.0$ - коэффициент потерь ценности, $c_e = 0.005$ - коэффициент энтропии.

\subsection{MaskablePPO}

Реализация использует библиотеку \texttt{sb3-contrib}, обеспечивающую интеграцию маски действий непосредственно в цикл PPO. При сборе роллаутов на каждом шаге запрашивается маска $m(\cdot, s_t)$; значение логитов невалидных действий устанавливается в $-\infty$ до вычисления softmax, что эквивалентно~\eqref{eq:masking}.

\subsection{Гиперпараметры}

Финальные гиперпараметры подобраны автоматическим перебором: 7 конфигураций обучались 100k шагов, победитель выбирался по средней награде на 100 оценочных эпизодах.

\begin{table}[H]
  \centering
  \caption{Гиперпараметры обучения}
  \label{tab:hparams}
  \begin{tabular}{lcc}
    \toprule
    \textbf{Параметр} & \textbf{MLP} & \textbf{CNN} \\
    \midrule
    \texttt{n\_steps}        & 512   & 512   \\
    \texttt{n\_epochs}       & 4     & 4     \\
    \texttt{batch\_size}     & 256   & 256   \\
    \texttt{learning\_rate}  & $10^{-4}$ & $7 \times 10^{-5}$ \\
    \texttt{ent\_coef}       & 0.005 & 0.005 \\
    \texttt{vf\_coef}        & 1.0   & 1.0   \\
    \texttt{clip\_range}     & 0.2   & 0.2   \\
    \texttt{gamma}           & 0.99  & 0.99  \\
    \texttt{gae\_lambda}     & 0.95  & 0.95  \\
    \texttt{n\_envs}         & 16    & 16    \\
    Всего шагов              & 5M    & 5M    \\
    \bottomrule
  \end{tabular}
\end{table}

% ════════════════════════════════════════════════════════════════
\section{Среда и инженерия признаков}
% ════════════════════════════════════════════════════════════════

\subsection{Реализация среды}

Среда реализована на основе стандарта \texttt{Gymnasium}~\cite{gymnasium2022} и предоставляет интерфейс \texttt{reset()} / \texttt{step(action)}, возвращающий тензор наблюдения, скалярную награду и флаг окончания эпизода. Для параллельного обучения используется \texttt{SubprocVecEnv} с 16 независимыми процессами; нормализация наблюдений выполняется через \texttt{VecNormalize}.

\subsection{Конструкция наблюдений: три поколения}

Ключевой исследовательский вопрос - \textit{как информационная насыщенность тензора состояния влияет на качество обучения?} Тензор наблюдений имеет форму $(C, 8, 8)$ и \\ эволюционировал через три поколения.

% ───────────────────────────────────────
\subsubsection{Поколение 1 - базовая архитектура $(C=6)$}
% ───────────────────────────────────────

\begin{table}[H]
  \centering
  \caption{Каналы наблюдения поколения 1}
  \label{tab:obs_gen1}
  \begin{tabular}{cl}
    \toprule
    \textbf{Канал} & \textbf{Содержимое} \\
    \midrule
    0 & Игровое поле $\in\{0,1\}^{8\times 8}$ (1 = занята) \\
    1 & Форма фигуры 1 в левом верхнем углу канала \\
    2 & Форма фигуры 2 в левом верхнем углу канала \\
    3 & Форма фигуры 3 в левом верхнем углу канала \\
    4 & Заполненность строк: $f_r[i]$ транслировано по строке $i$ \\
    5 & Заполненность столбцов: $f_c[j]$ транслировано по столбцу $j$ \\
    \bottomrule
  \end{tabular}
\end{table}

Каналы 4--5 несут \textit{декларативную} информацию о заполненности. Критик немедленно видит корреляцию $f_r = 0.875 \Rightarrow$ ``1 клетка до очистки $\Rightarrow$ высокая ценность'', не тратя тысячи шагов на вывод этой связи из бинарной доски.

% ───────────────────────────────────────
\subsubsection{Поколение 2 - карты размещения $(C=9)$}
% ───────────────────────────────────────

К каналам 0--5 добавлены три \textbf{placement heatmap}:
\begin{table}[H]
  \centering
  \caption{Дополнительные каналы поколения 2}
  \label{tab:obs_gen2}
  \begin{tabular}{cl}
    \toprule
    \textbf{Канал} & \textbf{Содержимое} \\
    \midrule
    6 & Карта размещения фигуры 1 \\
    7 & Карта размещения фигуры 2 \\
    8 & Карта размещения фигуры 3 \\
    \bottomrule
  \end{tabular}
\end{table}

Ячейка $(y,x)=1$ в канале $6+i$, если фигуру $i$ можно разместить с верхним левым углом в позиции $(x,y)$. Иначе говоря, эти каналы транслируют маску допустимых действий непосредственно в пространство наблюдений. Вычислительная стоимость нулевая: маска и так вычисляется для MaskablePPO.

Переход от \textit{декларативного} представления (``вот форма'') к \textit{процедурному} (``вот куда её можно поставить'') значительно снижает требования к аппроксимирующей способности нейросети.

% ───────────────────────────────────────
\subsubsection{Поколение 3 - топологическое зрение $(C=14)$}
% ───────────────────────────────────────

\begin{table}[H]
  \centering
  \caption{Дополнительные каналы поколения 3}
  \label{tab:obs_gen3}
  \begin{tabular}{cl}
    \toprule
    \textbf{Канал} & \textbf{Содержимое} \\
    \midrule
    9  & Survivability heatmap фигуры 1 \\
    10 & Survivability heatmap фигуры 2 \\
    11 & Survivability heatmap фигуры 3 \\
    12 & Largest Empty Blob Map \\
    13 & Dead Zone Mask \\
    \bottomrule
  \end{tabular}
\end{table}

\paragraph{Survivability heatmap.}
Для каждой допустимой позиции фигуры $i$ вычисляется, сколько суммарно допустимых позиций останется у двух оставшихся фигур после данного размещения, нормированное на $[0,1]$:
\begin{equation}
  \text{surv}_i(y,x)
  = \frac{\displaystyle\sum_{j \neq i} |\text{valid\_pos}(j,\, \mathcal{B} \oplus P_i^{(y,x)})|}{(N_{\text{fig}}-1)\cdot n^2},
  \label{eq:surv}
\end{equation}
где $\mathcal{B} \oplus P_i^{(y,x)}$ обозначает доску после размещения фигуры $i$ в позицию $(y,x)$, а $n=8$. Ячейка равна нулю, если фигуру нельзя разместить в данную позицию.

Канал~\eqref{eq:surv} обеспечивает \textbf{lookahead внутри раунда}: агент видит, насколько сохраняется свобода манёвра для оставшихся фигур, не выполняя явного дерева перебора.

\paragraph{Largest Empty Blob Map.}
Выполняется поиск в ширину (BFS) по пустым клеткам. Наибольшая связная компонента отмечается значением $1.0$, остальные - $0.0$. Агент получает визуальное представление ``основной территории'': изолированные мелкие фрагменты выделяются как потенциально мёртвые зоны.

\paragraph{Dead Zone Mask.}
BFS-компонента объявляется мёртвой, если ни одна из \textit{текущих трёх} фигур не может быть в неё размещена. Проверка выполняется именно против текущего раунда, а не против полного пула - иначе мономино 1×1 делало бы все компоненты живыми.

\subsection{Функция вознаграждения}

Разреженность вознаграждения устраняется с помощью плотного reward shaping~\cite{ng1999}. Итоговое вознаграждение на шаге $t$:
\begin{equation}
  r_t = r^{\text{place}} + r^{\text{lines}} + r^{\text{potential}} + r^{\text{round}} + r^{\text{death}},
\end{equation}
где компоненты определены в таблице~\ref{tab:reward}.

\begin{table}[H]
  \centering
  \caption{Компоненты функции вознаграждения}
  \label{tab:reward}
  \begin{tabular}{llr}
    \toprule
    \textbf{Компонент} & \textbf{Условие} & \textbf{Значение} \\
    \midrule
    $r^{\text{place}}$    & Размещение любой фигуры          & $+0.2$ \\
    $r^{\text{lines}}$    & За каждую очищенную линию         & $+1.0$ \\
    $r^{\text{potential}}$& За каждую размещённую клетку      & $f_r + f_c$ \\
    $r^{\text{round}}$    & Завершение раунда (Gen 3)         & $+0.5$ \\
    $r^{\text{death}}$    & Конец эпизода                     & $-1.0$ \\
    \bottomrule
  \end{tabular}
\end{table}

\paragraph{Cell potential.}
Ключевой компонент --- \textit{cell\_potential}: каждая размещённая клетка приносит вознаграждение, пропорциональное заполненности соответствующей строки $f_r[i]$ и столбца $f_c[j]$:
\begin{equation}
  r^{\text{potential}} = \alpha \sum_{(i,j) \in P} \bigl( f_r[i] + f_c[j] \bigr),
  \qquad \alpha = 0.1,
\end{equation}
где $P$ --- множество занятых фигурой клеток. Коэффициент $\alpha = 0.1$ подобран так, чтобы вклад \textit{cell\_potential} за один ход (${\approx}0.1$--$0.8$) был соизмерим с $r^{\text{place}} = 0.2$, но не перекрывал сигнал $r^{\text{lines}} = 1.0$ от реальной очистки. Это создаёт замкнутую обратную связь с каналами 4--5 наблюдения: критик видит заполненность строк и столбцов, а cell\_potential вознаграждает за её использование.

Компонент $r^{\text{round}}=+0.5$, введённый в поколении 3, явно поощряет ``дожить до следующей тройки'', что согласовано с выживаемостным сигналом survivability heatmap.

\subsection{Оптимизация производительности среды}

Скорость среды критична при обучении. Исходная реализация на чистом Python выполняла ${\approx}253\,000$ вызовов \texttt{can\_place} за один шаг через вложенные циклы. Последовательная оптимизация снизила время шага с $5.6$~мс до $0.27$~мс (ускорение ${\approx}20\times$):

\begin{table}[H]
  \centering
  \caption{Этапы оптимизации среды ($8\times8$)}
  \label{tab:perf}
  \small
  \begin{tabular}{lrr}
    \toprule
    \textbf{Этап} & \textbf{8×8, мс/шаг} & \textbf{20×20, мс/шаг} \\
    \midrule
    Исходный код (Python-циклы)               & $5.6$ & ${\sim}8{+}$ \\
    + \texttt{correlate2d} (scipy)             & $3.5$ & $3.4$ \\
    + Слияние двух BFS-проходов в один         & $0.47$ & ${\sim}1.3$ \\
    + Cython (\texttt{\_fast.pyx})             & ${\sim}0.5$ & ${\sim}0.5$ \\
    + FFT-кэш + дедупликация $A_j$             & $\mathbf{0.27}$ & $\mathbf{0.23}$ \\
    \bottomrule
  \end{tabular}
\end{table}

Наибольший скачок даёт слияние BFS: \texttt{\_largest\_empty\_blob} и \texttt{\_dead\_zone\_mask} обходили одни и те же пустые клетки двумя независимыми проходами --- объединение в один проход даёт $7\times$ ускорение шага. FFT-кэш устраняет повторное вычисление $\mathrm{FFT}(B_\text{ext})$ в survivability heatmap: тот зависит только от пары фигур $(i, j)$, а не от состояния доски, поэтому кэшируется по ключу $(i_{\text{pool}}, j_{\text{pool}})$.

% ════════════════════════════════════════════════════════════════
\section{Архитектуры нейронных сетей}
% ════════════════════════════════════════════════════════════════

\subsection{MLP-агент}

MLP-агент разворачивает тензор наблюдений $(C, 8, 8)$ в плоский вектор размерности $C \times 64$ и обрабатывает его \textit{раздельными} сетями актора и критика:
\begin{align}
  &\text{Актор: } [C\cdot 64] \to [256] \to [256] \to [192], \\
  &\text{Критик: } [C\cdot 64] \to [512] \to [512] \to [1].
\end{align}

Раздельные сети важны: единая общая сеть создаёт конкуренцию градиентов между политикой и критиком. MLP видит весь вектор состояния одновременно, без индуктивного смещения на пространственную локальность - что оказывается преимуществом на маленьком поле 8×8, где глобальная картина важнее локальных паттернов.

\subsection{CNN-агент (SmallCNN)}

Стандартный \texttt{CnnPolicy} библиотеки Stable Baselines 3 использует NatureCNN~\cite{mnih2015} с шагом свёртки $\mathrm{stride}=4$ и $\mathrm{stride}=2$. На поле $8 \times 8$ такие страйды схлопывают пространственное разрешение до нуля, вызывая \texttt{RuntimeError}. Разработан кастомный экстрактор \textbf{SmallCNN}:

\begin{equation}
  (C, 8, 8)
  \xrightarrow{\text{Conv}_{3\times3}^{32}, \text{ReLU}}
  (32, 8, 8)
  \xrightarrow{\text{Conv}_{3\times3}^{64}, \text{ReLU}}
  (64, 8, 8)
  \xrightarrow{\text{Conv}_{3\times3}^{64}, \text{ReLU}}
  (64, 8, 8)
  \xrightarrow{\text{Flatten}}
  [4096]
  \xrightarrow{\text{Linear}}
  [256].
\end{equation}

Все свёрточные слои используют $\mathrm{stride}=1$, $\mathrm{padding}=1$ - пространственный размер $8\times 8$ сохраняется через все три слоя. Рецептивное поле финального слоя составляет $7\times 7$, то есть охватывает почти всю доску из любой позиции.

\paragraph{Ключевые технические решения.}
\begin{enumerate}
  \item \textbf{BatchNorm исключён намеренно.} На этапе тестирования была добавлена батч-нормализация по аналогии с Computer Vision. Результат: \texttt{clip\_fraction} вырос до 0.45, \texttt{explained\_variance} упал до $-0.57$. Причина: в RL-цикле модель переключается между \texttt{eval()} при сборе роллаутов и \texttt{train()} при обновлении весов. BatchNorm имеет различное поведение в этих режимах, что искажает оценки advantage и разрушает критика.
  \item \textbf{Автоматический размер Linear.} Вместо хардкода числа 4096 используется форвард-проход через фиктивный тензор при инициализации, что делает экстрактор инвариантным к размеру поля.
  \item \textbf{Чувствительность к learning rate.} При $\eta = 3\times 10^{-5}$ обновления политики почти отсутствовали (\texttt{clip\_fraction} = 0.002--0.005), при $\eta = 7\times 10^{-5}$ обучение нормализовалось.
\end{enumerate}

\subsection{HierarchicalCNN (поколение 4)}

Для поля $16\times16$ разработан \textbf{HierarchicalCNN} с иерархическим понижением разрешения:

\begin{equation}
	\begin{aligned}
		(15,16,16)
		&\xrightarrow{\text{Conv}^{32},\text{BN,ReLU}} (32,16,16) \\
		&\xrightarrow{\text{Conv}^{64},\text{stride=2}} (64,8,8) \\
		&\xrightarrow{\text{Conv}^{128},\text{stride=2}} (128,4,4) \\
		&\xrightarrow{\text{Conv}^{256},\text{stride=2}} (256,2,2) \\
		&\xrightarrow{\text{Flatten}\to\text{Linear}} [512].
	\end{aligned}
\end{equation}

Все свёрточные слои имеют ядро $3\times3$, padding$=1$, и сопровождаются BatchNorm и ReLU. Число параметров --- $1\,292\,800$ (${\approx}1.29$М). Головы актора и критика: $\pi=[512,256]$, $V=[1024,512]$.

В отличие от SmallCNN (поколения 1--3), HierarchicalCNN содержит BatchNorm. Это допустимо благодаря BC warm-start: политика стартует достаточно хорошо, чтобы статистики BN были устойчивы с первых шагов PPO. Дополнительно VecNormalize нормализует наблюдения на уровне окружения, снижая чувствительность к переключению eval/train.

\subsection{SmallViT (поколение 4)}

\textbf{SmallViT} применяет трансформерную архитектуру~\cite{dosovitskiy2020} к пространственным наблюдениям:

\begin{enumerate}
  \item \textbf{Patch embedding.} Поле $16\times16$ разбивается на $16$ патчей $4\times4$. Каждый патч разворачивается в вектор размерности $15\times4\times4=240$ и проецируется в $\mathbb{R}^{256}$ линейным слоем.
  \item \textbf{CLS-токен и позиционные эмбеддинги.} К последовательности $16$ патчей добавляется обучаемый CLS-токен; все $17$ позиций получают обучаемые позиционные эмбеддинги.
  \item \textbf{TransformerEncoder.} $2$ слоя с Pre-LN (\texttt{norm\_first=True}), каждый --- мульти-головое само-внимание (embed\_dim$=256$, $n_\text{heads}=4$) и двуслойный FFN ($256\to512\to256$).
  \item \textbf{Агрегация.} Патчи усредняются, результат нормализуется LayerNorm и проецируется $\mathrm{Linear}(256\to512)$ c ReLU.
\end{enumerate}

Число параметров --- $1\,252\,096$ (${\approx}1.25$М) --- практически совпадает с HierarchicalCNN, что обеспечивает \textit{честное} архитектурное сравнение при равном параметрическом бюджете.

\paragraph{Выбор гиперпараметров.}
Размер патча $4\times4$ выбран как компромисс: меньший патч ($2\times2$) даёт $64$ токена --- слишком длинная последовательность при embed\_dim$=256$ при данном параметрическом бюджете; больший патч ($8\times8$) --- лишь $4$ токена, недостаточно для захвата пространственной структуры. Два слоя трансформера достаточны для агрегации информации между всеми $16$ патчами за счёт глобального внимания с первого слоя. Четыре головы внимания ($\text{head\_dim}=64$) соответствуют эмпирической рекомендации head\_dim $\in [32, 64]$ для малых трансформеров~\cite{dosovitskiy2020}.

Механизм само-внимания позволяет любому патчу напрямую взаимодействовать с любым другим вне зависимости от расстояния на доске. Это индуктивное смещение более подходит для задачи, где пространственные зависимости нелокальны (размещение фигуры в одном углу влияет на выживаемость позиции в другом).

% ════════════════════════════════════════════════════════════════
\section{Бейзлайны}
% ════════════════════════════════════════════════════════════════

\subsection{RandomAgent}

Случайный агент равномерно выбирает ход из множества допустимых действий на каждом шаге. Обеспечивает нижнюю границу производительности.

\subsection{HeuristicAgent}

Жадный алгоритм, вдохновлённый методом Dellacherie для игры Tetris~\cite{dellacherie2003}. На каждом ходу симулируются все допустимые размещения, и результирующая доска $\mathcal{B}'$ оценивается по взвешенной сумме семи признаков:
\begin{equation}
  \text{score}(\mathcal{B}') = \sum_{k=1}^{7} w_k\, \phi_k(\mathcal{B}').
\end{equation}

\begin{table}[H]
  \centering
  \caption{Признаки и веса HeuristicAgent}
  \label{tab:heuristic}
  \small
  \begin{tabular}{clrl}
    \toprule
    $k$ & \textbf{Признак} $\phi_k$ & \textbf{Вес} $w_k$ & \textbf{Описание} \\
    \midrule
    1 & \texttt{lines\_cleared}  & $+10.0$ & Число очищенных строк + столбцов \\
    2 & \texttt{combo\_bonus}    & $+5.0$  & Бонус $\propto L(L{-}1)/2$ за $L\geq2$ линий за раз \\
    3 & \texttt{perfect\_clear}  & $+50.0$ & Бонус за полное очищение доски \\
    4 & \texttt{holes}           & $-6.0$  & Пустые клетки, не достижимые BFS от границы \\
    5 & \texttt{near\_complete}  & $+3.0$  & Число строк/столбцов, заполненных на $\geq7/8$ \\
    6 & \texttt{adjacency}       & $+0.4$  & Число пар соседних занятых клеток (компактность) \\
    7 & \texttt{fill\_variance}  & $-0.6$  & Дисперсия заполнения строк и столбцов \\
    \bottomrule
  \end{tabular}
\end{table}

Признак \textit{holes} (вес $-6$) является наиболее важным ограничительным: замурованные клетки не могут быть очищены и необратимо «съедают'' пространство доски. Именно этот признак послужил прообразом канала~14 тензора наблюдений Gen~4 (Holes-from-border mask). Выбирается ход, максимизирующий оценку.

\begin{remark}
  Вычислительная сложность эвристики составляет $O(|\text{valid}| \times n^2)$ на ход, где $|\text{valid}| \approx 180$ и $n=8$, итого $\approx 11\,500$ операций. На поле $16\times 16$ это возрастает в 16 раз. RL-агент выполняет \textbf{один форвард} нейронной сети, не зависит от размера поля.
\end{remark}

% ════════════════════════════════════════════════════════════════
\section{Эксперименты и результаты}
% ════════════════════════════════════════════════════════════════

Все агенты оценивались на 500 независимых эпизодах с фиксированным \texttt{seed=42}. Метрики:
\begin{itemize}
  \item \textbf{Lines cleared} - число очищенных линий за эпизод;
  \item \textbf{Reward} - суммарная нескалярная награда;
  \item \textbf{Pieces placed} - число размещённых фигур;
  \item \textbf{CV} = $\sigma / \mu$ - коэффициент вариации (мера стабильности).
\end{itemize}

\subsection{Сводные результаты}

\begin{table}[H]
  \centering
  \caption{Результаты оценки агентов (500 эпизодов, seed=42)}
  \label{tab:results_all}
  \small
  \begin{tabular}{lrrr}
    \toprule
    \textbf{Агент} & \textbf{Reward} $\mu\pm\sigma$ & \textbf{Lines} $\mu\pm\sigma$ & \textbf{Pieces} $\mu\pm\sigma$ \\
    \midrule
    Random            & $10.48 \pm 11.67$ & $1.79 \pm 2.12$  & $16.32 \pm 5.37$  \\
    Heuristic         & $45.91 \pm 40.10$ & $8.48 \pm 7.60$  & $32.31 \pm 17.73$ \\
    \midrule
    MLP Gen 1         & $48.23 \pm 31.45$ & $8.84 \pm 5.90$  & $33.43 \pm 13.95$ \\
    CNN Gen 1         & $39.92 \pm 25.23$ & $7.36 \pm 4.74$  & $29.40 \pm 11.26$ \\
    \midrule
    MLP Gen 2         & $53.89 \pm 35.13$ & $9.88 \pm 6.57$  & $35.85 \pm 15.63$ \\
    CNN Gen 2         & $46.36 \pm 30.50$ & $8.50 \pm 5.72$  & $32.68 \pm 13.58$ \\
    \midrule
    MLP Gen 3         & $\mathbf{60.76 \pm 40.96}$ & $\mathbf{11.21 \pm 7.68}$ & $\mathbf{38.81 \pm 18.14}$ \\
    CNN Gen 3 (нуля)  & $26.63 \pm 20.44$ & $4.73 \pm 3.79$  & $23.63 \pm 9.03$  \\
    CNN Gen 3 (TL)    & $48.23 \pm 30.17$ & $8.85 \pm 5.66$  & $33.52 \pm 13.45$ \\
    \bottomrule
  \end{tabular}
\end{table}

\subsection{Поколение 1}

MLP Gen 1 превзошёл эвристику по среднему числу линий на $+4.2\%$ ($8.84$ vs $8.48$). Оба нейросетевых агента существенно стабильнее эвристики: $\mathrm{CV}_\text{MLP} = 0.67$, $\mathrm{CV}_\text{CNN} = 0.64$ против $\mathrm{CV}_\text{heur} = 0.90$ - преимущество $-28\%$.

Стабильность является принципиальным преимуществом: эвристика нестабильна по своей природе - одношаговый жадный выбор без учёта будущих фигур даёт $20+$ линий при удачных раскладах и лишь $3$--$5$ при сложных.

\subsection{Поколение 2}

Добавление placement heatmap (каналы 6--8) дало значительный прирост для MLP:
\begin{equation}
  \Delta\text{Lines}^{\text{MLP}}_{1\to2} = \frac{9.88 - 8.84}{8.84} \approx +11.8\%.
\end{equation}

MLP Gen 2 уверенно превзошёл эвристику на $+16.5\%$ по lines и $+17\%$ по reward. CNN при оптимальном $\eta = 7\times10^{-5}$ вышла на паритет с эвристикой ($8.50$ vs $8.48$).

\subsection{Поколение 3 и проблема порочного круга}

Обучение CNN с нуля на тензоре $(14, 8, 8)$ привело к регрессии до $4.73$ линий. Анализ диагностических метрик позволил установить причину:
\begin{equation}
  \underbrace{\text{EV}=0.25}_{\text{слабый критик}}
  \;\Rightarrow\;
  \text{шумный }\hat{A}_t
  \;\Rightarrow\;
  \text{нестабильная политика}
  \;\Rightarrow\;
  \text{шумные topology-каналы}
  \;\Rightarrow\;
  \underbrace{\text{EV}=0.25}_{\circlearrowleft}
\end{equation}

Topology-каналы (survivability, blob, dead zone) информативны \textit{только при достаточно хорошей политике}. Пока политика случайна, survivability представляет собой шум - критик не может на него опереться, и порочный круг замыкается.

Динамика кривой Gen 3 (с нуля) на рис.~\ref{fig:learning_cnn_gen123} объясняется неоднородностью новых каналов по степени зависимости от качества политики. Среди пяти дополнительных каналов часть является \textbf{детерминированными} --- dead zones (клетки, недостижимые при текущих фигурах) и blob-связность вычисляются непосредственно из состояния доски и несут полезную геометрическую информацию независимо от того, насколько хороша политика. Именно они объясняют начальный рост Gen 3, превышающий темп Gen 2: агент получает больше значимых признаков с первых шагов.

Однако канал \textbf{survivability} --- оценка числа доступных будущих ходов --- зависит от качества самой политики. Пока политика случайна, survivability отражает лишь случайный исход, а не реальную «выживаемость» позиции. По мере того как CNN начинает опираться на этот канал, оценки критика дестабилизируются, advantage $\hat{A}_t$ становится ненадёжным, и политика деградирует --- отсюда спад после $\approx$3М шагов. Таким образом, наличие policy-зависимых каналов при обучении с нуля создаёт петлю обратной связи, которой нет ни у MLP (плоская обработка), ни у Gen 3 с трансфером (старт с рабочей политики, где survivability уже информативен).

MLP не страдает от этой проблемы: он обрабатывает тензор как плоский вектор и способен выучить малые веса для новых шумных каналов, не нарушая уже выученных весов для базовых.

\subsection{Transfer Learning для CNN}
\label{sec:transfer}

Для разрыва порочного круга применён \textbf{transfer learning}: веса SmallCNN из Gen 2 (9 каналов) переносятся в Gen 3 (14 каналов) с нулевой инициализацией новых каналов:

\begin{align}
  \text{Conv1}^{(14)} &= \begin{bmatrix}
    W^{(9)}_{\text{Gen2}} \;\big|\; \mathbf{0}_{32\times5\times3\times3}
  \end{bmatrix}, \\
  \text{Conv2, Conv3, Linear, головы} &= \text{скопированы без изменений}.
\end{align}

При fine-tuning: $\eta = 5\times10^{-5}$, $\|\nabla\|_{\max} = 0.3$ (grad clipping). Агент стартует с уже рабочей геометрической базой ($\approx\!8.7$ lines), survivability-каналы с первых шагов несут информативный сигнал - порочный круг не возникает.

Данный подход фактически реализует \textbf{неявный curriculum learning}: базовые навыки геометрического размещения освоены заранее, топологические признаки добавляются поверх рабочей политики.

Результат: CNN Gen 3 (TL) достигает $8.85$ lines против $8.50$ у Gen 2 - небольшой прирост, но принципиально важен сам факт: transfer learning спасает CNN там, где обучение с нуля провалилось.

\begin{figure}[H]
  \centering
  \includegraphics[width=0.9\textwidth]{../results/charts/report/fig_learning_cnn_gen123.png}
  \caption{Кривые обучения CNN через три поколения. CNN Gen~3 (с нуля) деградирует к нулю --- проблема порочного круга. CNN Gen~3 (трансфер) стартует с уже рабочей политики и продолжает рост. Пунктир --- уровень HeuristicAgent.}
  \label{fig:learning_cnn_gen123}
\end{figure}

\subsection{Сравнительная таблица поколений}

\begin{table}[H]
  \centering
  \caption{Эволюция пространства наблюдений и результатов}
  \label{tab:gen_compare}
  \small
  \begin{tabular}{lccc}
    \toprule
    & \textbf{Gen 1} & \textbf{Gen 2} & \textbf{Gen 3 (TL)} \\
    \midrule
    Форма тензора       & $(6,8,8)$ & $(9,8,8)$ & $(14,8,8)$ \\
    Тип информации      & Декларативная & + Процедурная & + Lookahead \\
    Round bonus         & --- & --- & $+0.5$ \\
    Lines MLP (5M)      & $8.84$ & $9.88$ (+12\%) & $11.21$ (+27\%) \\
    Lines CNN (5M)      & $7.36$ & $8.50$ (+15\%) & $8.85$ (+4\%) \\
    CV CNN              & $0.64$ & $0.67$ & $0.64$ \\
    MLP лучше эвристики & $+4\%$ & $+16.5\%$ & $+32\%$ \\
    \bottomrule
  \end{tabular}
\end{table}

% ════════════════════════════════════════════════════════════════
\section{Поколение 4 --- Behavior Cloning и масштабирование до \texorpdfstring{$16\times16$}{16x16}}
% ════════════════════════════════════════════════════════════════

\subsection{Мотивация масштабирования}

На поле $8\times8$ агент Gen~3 превзошёл эвристику на $4$--$32\%$ в зависимости от архитектуры. Гипотеза: при увеличении поля до $16\times16$ преимущество нейросети должно существенно возрасти по трём причинам.

Во-первых, эвристика остаётся одношаговой жадной: на большой доске цена ошибки выше, мёртвые зоны крупнее, а паттерны сложнее. Во-вторых, RL-агент оптимизирует глобальный сигнал и способен выработать нелокальные стратегии. В-третьих, вычислительная сложность эвристики на одном ходу составляет $O(|\text{valid}| \times n^2)$, где $|\text{valid}| \sim n^2$ --- число допустимых позиций. При увеличении поля оба множителя растут, итого $O(n^4)$; тогда как один форвард нейросети не зависит от $n$.

Результаты Gen~4 подтверждают гипотезу: CNN (BC+PPO) набирает \textbf{18.76 lines} против \textbf{7.21 lines} у эвристики --- разрыв $\mathbf{2.6\times}$$.

\begin{remark}
  На поле $16\times16$ ситуация принципиально иная по сравнению с $8\times8$. Рецептивное поле SmallCNN ($7\times7$) покрывало бы лишь ${\approx}19\%$ доски $16\times16$, поэтому в Gen~4 он заменён на HierarchicalCNN с иерархическим понижением разрешения. Пространство действий вырастает до $768$; входной вектор MLP расширялся бы до $15\times256=3840$ элементов, тогда как HierarchicalCNN сохраняет компактный $512$-мерный эмбеддинг независимо от размера поля --- именно поэтому на $16\times16$ CNN оказывается более пригодной архитектурой, чем MLP.
\end{remark}

\subsection{Проблема холодного старта}

На $16\times16$ пространство действий составляет $|\mathcal{A}|=3\times16\times16=768$. Случайный агент набирает лишь ${\approx}0.5$ линии/эп. Прямое PPO-обучение со случайной инициализацией застревает ниже $1$ line/эп. в течение первого $1$М шагов: политика не получает значимого обучающего сигнала и не выходит из области деструктивного исследования.

Корень проблемы --- разреженность сигнала. На $8\times8$ случайный агент очищает ${\approx}1.8$ линии/эп., что даёт достаточную обратную связь для старта PPO. На $16\times16$ та же случайная политика производит ${\approx}0.5$ линии/эп. --- сигнал слишком слаб для корректного обновления advantage.

\subsection{Behavior Cloning как тёплый старт}

\textbf{Behavior Cloning} (BC)~\cite{pomerleau1991} --- метод имитационного обучения, в котором политика обучается воспроизводить действия эксперта через минимизацию потерь кросс-энтропии. Это supervised learning, не требующий взаимодействия со средой на этапе предобучения:
\begin{equation}
  \mathcal{L}_{\text{BC}}(\theta)
  = -\frac{1}{N}\sum_{i=1}^{N} \log \pi_\theta\!\left(a^{(i)}_{\text{exp}} \mid s^{(i)}\right),
  \label{eq:bc}
\end{equation}
где $\bigl\{(s^{(i)},\, a^{(i)}_{\text{exp}})\bigr\}_{i=1}^{N}$ --- датасет из $N=500\,000$ пар состояние--действие, собранных от HeuristicAgent.

\paragraph{Двухфазный пайплайн.}

\begin{enumerate}
  \item \textbf{Фаза 1 --- Behavior Cloning.} Собирается датасет из $500\,000$ пар $(s, a_{\text{exp}})$ путём прогонки HeuristicAgent. Нейронная сеть обучается минимизировать: $50$ эпох, batch\_size$=512$, $\eta=10^{-3}$ с cosine annealing. Итоговая top-1 точность: $40$--$45\%$ (доля шагов, на которых $\arg\max\,\pi_\theta(s) = a_{\text{exp}}$). Низкое значение объясняется стохастичностью эвристики при равных score: среди нескольких геометрически равноценных ходов эвристика выбирает случайный, поэтому нейронная сеть воспроизводит верную геометрию, но не угадывает конкретный тайbreak. После фазы~1 политика стартует с ${\approx}4$ lines/эп.

  \item \textbf{Фаза 2 --- PPO fine-tuning.} Веса актора загружаются из BC-чекпоинта; критик \\ инициализируется случайно. Проводится $15$М шагов MaskablePPO c VecNormalize (\texttt{norm\_reward=True}, clip$=10$) и линейным расписанием скорости обучения $\eta: 5\times10^{-5}\to5\times10^{-6}$. Обе фазы выполнялись на Kaggle T4 GPU: фаза~1 заняла ${\approx}1$~ч., фаза~2 --- ${\approx}7$~ч.
\end{enumerate}

Принципиально важно, что критик \textit{не} переносится из BC: BC-политика предоставляет хорошую начальную точку, а PPO строит корректную оценку функции ценности заново, исключая систематическое смещение «экспертного» критика, который не оптимизировался по RL-объективу.

Подход аналогичен описанному в разделе~\ref{sec:transfer} curriculum learning: базовые навыки передаются supervised-образом, а тонкая настройка под глобальный reward выполняется PPO.

\subsection{Пространство наблюдений $(15, 16, 16)$}

Каналы $0$--$13$ масштабированы на поле $16\times16$. Добавлен канал~14:

\begin{table}[H]
  \centering
  \caption{Канал 14 поколения 4}
  \label{tab:obs_gen4}
  \begin{tabular}{cl}
    \toprule
    \textbf{Канал} & \textbf{Содержимое} \\
    \midrule
    14 & \textit{Holes-from-border mask}: пустые клетки, недостижимые от границы доски (BFS) \\
    \bottomrule
  \end{tabular}
\end{table}

В отличие от Dead Zone (канал~13, зависящего от текущих фигур), \textit{Holes-from-border} --- топологический инвариант доски: замурованные клетки, которые не могут быть очищены независимо от будущих фигур. Канал соответствует признаку \texttt{holes} (с весом $-6$) в эвристике, но передаётся критику в виде пространственной маски, обеспечивая явное понимание структурных потерь позиции.

\subsection{Результаты Gen 4}

Оценка проводилась на $500$ независимых эпизодах с \texttt{seed=42} на поле $16\times16$.

\begin{table}[H]
  \centering
  \caption{Результаты на поле $16\times16$ (500 эпизодов, seed=42)}
  \label{tab:results_gen4}
  \begin{tabular}{lrrrr}
    \toprule
    \textbf{Агент} & \textbf{Lines} $\mu$ & \textbf{Lines} $\sigma$ & \textbf{Reward} $\mu$ & \textbf{Pieces} $\mu$ \\
    \midrule
    Random           & $0.47$  & $0.89$  & $16.29$  & $56.83$  \\
    Heuristic        & $7.21$  & $7.09$  & $64.23$  & $94.17$  \\
    \midrule
    CNN (BC+PPO)     & $\mathbf{18.76}$ & $9.96$ & $\mathbf{123.23}$ & $\mathbf{148.25}$ \\
    ViT (BC+PPO)     & $18.12$ & $\mathbf{9.74}$  & $121.76$ & $147.01$ \\
    \bottomrule
  \end{tabular}
\end{table}

Оба нейросетевых агента превзошли эвристику в $2.6\times$ по числу очищенных линий. Разница между CNN и ViT составляет $0.64$ lines $\approx 0.06\sigma$ и статистически незначима (см.~раздел~\ref{sec:stat_gen4}).

\begin{figure}[H]
  \centering
  \includegraphics[width=0.75\textwidth]{../results/charts/report/fig_comparison_gen4.png}
  \caption{Среднее число очищенных линий ($\pm\sigma$) для каждого агента на доске $16\times16$ ($n=500$ эпизодов).}
  \label{fig:comparison_gen4}
\end{figure}

\begin{figure}[H]
  \centering
  \includegraphics[width=0.9\textwidth]{../results/charts/report/fig_learning_gen4.png}
  \caption{Кривые обучения Gen~4 в фазе PPO fine-tuning. По горизонтали --- число шагов взаимодействия со средой; по вертикали --- скользящее среднее награды за эпизод. Пунктир --- уровень HeuristicAgent.}
  \label{fig:learning_gen4}
\end{figure}

\subsection{Наблюдение: диагональная стратегия}

Все нейросетевые агенты заполняют поле по диагонали из левого верхнего угла. С точки зрения классического Тетриса это нелогично, однако для Block Blast оптимально: в этой игре одновременно очищаются строки \textit{и} столбцы. Диагональное заполнение одновременно прогрессирует оба направления, максимизируя вероятность combo-очистки. Нейросеть выработала эту стратегию самостоятельно через reward shaping --- без каких-либо явных указаний.

На рис.~\ref{fig:diagonal_board} приведено состояние доски на шаге~30 одного эпизода для обоих нейросетевых агентов. Несмотря на различные архитектуры, CNN и ViT демонстрируют идентичный паттерн: заполнение идёт диагональной волной из левого верхнего угла.

\begin{figure}[H]
  \centering
  \includegraphics[width=1.05\textwidth]{../results/charts/report/fig_diagonal_step30.png}
  \caption{Состояние доски $16\times16$ на шаге~30 одного эпизода. Слева --- CNN (BC+PPO), справа --- ViT (BC+PPO). Оба нейросетевых агента демонстрируют одинаковую диагональную стратегию заполнения из левого верхнего угла.}
  \label{fig:diagonal_board}
\end{figure}

% ════════════════════════════════════════════════════════════════
\section{Статистический анализ}
% ════════════════════════════════════════════════════════════════

Статистический анализ выполнен на 500 независимых эпизодах. Применялись:
\begin{itemize}
  \item \textbf{Критерий Манна--Уитни} (непараметрический, без предположения о нормальности);
  \item \textbf{CLES} (Common Language Effect Size) - $P(\text{агент} > \text{эвристика})$;
  \item \textbf{Cohen's d} - стандартизированный размер эффекта;
  \item \textbf{Bootstrap 95\% CI} для разности средних (1000 ресемплов);
  \item \textbf{Критерий Вилкоксона со знаковыми рангами} (парные тесты на одинаковых начальных фигурах).
\end{itemize}

\subsection{MLP Gen 3 против эвристики}

\begin{table}[H]
  \centering
  \caption{Статистические тесты: MLP Gen 3 vs Heuristic (Lines cleared)}
  \label{tab:stat_mlp3}
  \begin{tabular}{lr}
    \toprule
    \textbf{Метрика} & \textbf{Значение} \\
    \midrule
    Mann--Whitney $U$ (p-value) & $p < 0.001$ \\
    CLES & $0.634$ \\
    Cohen's $d$ & $+0.348$ \\
    Bootstrap 95\% CI для $\Delta\mu$ & $[+2.02;\; +3.37]$ \\
    Win rate (парный тест) & $59.4\%$ \\
    Медианная разность (парный) & $+2.00$ линии \\
    \bottomrule
  \end{tabular}
\end{table}

Значение $\mathrm{CLES}=0.634$ означает, что в $63.4\%$ случайно выбранных пар эпизодов MLP Gen 3 показывает результат лучше эвристики. Cohen's $d = 0.348$ соответствует малому--среднему размеру эффекта ($d=0.2$ - мало, $d=0.5$ - средне).

\subsection{Таблица выживаемости}

\begin{table}[H]
  \centering
  \caption{Доля эпизодов, достигших $\geq N$ очищенных линий}
  \label{tab:survival}
  \begin{tabular}{lrrrrrrr}
    \toprule
    \textbf{Агент} & $\geq1$ & $\geq5$ & $\geq10$ & $\geq15$ & $\geq20$ & $\geq25$ & $\geq30$ \\
    \midrule
    Random       & 66\% & 12\% &  2\% &  0\% &  0\% &  0\% &  0\% \\
    Heuristic    & 94\% & 64\% & 33\% & 17\% &  9\% &  4\% &  3\% \\
    MLP Gen 3    & 99\% & \textbf{85\%} & \textbf{49\%} & \textbf{26\%} & \textbf{14\%} &  \textbf{7\%} &  \textbf{3\%} \\
    CNN Gen 3 TL & 99\% & 78\% & 40\% & 15\% &  6\% &  2\% &  1\% \\
    \bottomrule
  \end{tabular}
\end{table}

Доля эпизодов с $\geq 10$ линиями у MLP Gen 3 составляет $49\%$ против $33\%$ у эвристики --- относительный прирост $+48\%$.

\subsection{Gen 4: CNN и ViT против эвристики (\texorpdfstring{$16\times16$}{16x16})}
\label{sec:stat_gen4}

Статистические тесты проведены на $500$ эпизодах при \texttt{seed=42}. Критерий Манна--Уитни выбран вместо $t$-теста ввиду тяжёлых хвостов распределения числа линий. CLES, Cohen's~$d$ и Bootstrap~CI вычислены дополнительно.

\begin{table}[H]
  \centering
  \caption{Статистические тесты Gen 4 vs Heuristic (lines cleared, $16\times16$)}
  \label{tab:stat_gen4}
  \small
  \begin{tabular}{lrr}
    \toprule
    \textbf{Метрика} & \textbf{CNN vs Heuristic} & \textbf{ViT vs Heuristic} \\
    \midrule
    Mann--Whitney $p$-value              & $<0.001$                  & $<0.001$                  \\
    CLES                                 & $0.852$                   & $0.847$                   \\
    Cohen's $d$                          & $+1.39$                   & $+1.34$                   \\
    Bootstrap 95\% CI ($\Delta\mu$)      & $[{+10.62};\;{+12.69}]$   & $[{+10.11};\;{+12.18}]$   \\
    Paired win rate                      & $84.2\%$                  & $85.8\%$                  \\
    \bottomrule
  \end{tabular}
\end{table}

Cohen's $d = 1.39$ соответствует очень большому размеру эффекта (порог по Коэну: $d > 0.8$ --- «большой»). CLES $= 0.852$ означает, что CNN выигрывает в $85\%$ произвольно выбранных пар эпизодов. Bootstrap CI по разности средних не включает нуль во всех сравнениях --- эффект статистически достоверен и воспроизводим.

\paragraph{CNN против ViT.} Разность средних составляет $0.64$ lines при $\sigma \approx 10$ --- оба агента статистически неразличимы ($p = 0.42$, Mann--Whitney U). Это указывает на то, что ограничивающим фактором является качество BC+PPO пайплайна, а не выбор архитектуры (CNN vs Transformer).

\subsection{Таблица выживаемости Gen 4}

\begin{table}[H]
  \centering
  \caption{Доля эпизодов, достигших $\geq N$ очищенных линий ($16\times16$, 500 эп.)}
  \label{tab:survival_gen4}
  \begin{tabular}{lrrrrrrr}
    \toprule
    \textbf{Агент} & $\geq1$ & $\geq5$ & $\geq10$ & $\geq15$ & $\geq20$ & $\geq25$ & $\geq30$ \\
    \midrule
    Random          & $29\%$  & $1\%$   & $0\%$  & $0\%$  & $0\%$  & $0\%$  & $0\%$  \\
    Heuristic       & $88\%$  & $51\%$  & $32\%$ & $15\%$ & $7\%$  & $3\%$  & $1\%$  \\
    CNN (BC+PPO)    & $100\%$ & $97\%$  & $85\%$ & $63\%$ & $41\%$ & $27\%$ & $13\%$ \\
    ViT (BC+PPO)    & $100\%$ & $97\%$  & $84\%$ & $61\%$ & $41\%$ & $23\%$ & $13\%$ \\
    \bottomrule
  \end{tabular}
\end{table}

На пороге $\geq 30$ lines нейросети превосходят эвристику в $13\times$ ($13\%$ против $1\%$). Начиная с $\geq 1$ line оба нейросетевых агента достигают $100\%$ --- то есть никогда не завершают эпизод без единой очищенной линии, тогда как эвристика в $12\%$ случаев не набирает ни одной.

\begin{figure}[H]
  \centering
  \includegraphics[width=1\textwidth]{../results/charts/report/fig_survival_gen4.png}
  \caption{Survival Analysis Gen~4: доля эпизодов, достигших заданного числа очищенных линий. Нейросетевые агенты сохраняют значительное преимущество над эвристикой на всех порогах.}
  \label{fig:survival_gen4}
\end{figure}

% ════════════════════════════════════════════════════════════════
\section{Заключение}
% ════════════════════════════════════════════════════════════════

В работе исследована эволюция RL-агента для игры Block Blast! через четыре поколения, в каждом из которых систематически наращивались информационное содержание тензора наблюдений, сложность архитектуры и масштаб задачи.

\paragraph{Основные результаты.}

\begin{enumerate}
  \item \textbf{Инженерия наблюдений --- ключевой фактор качества.} Добавление процедурных признаков (placement heatmaps, Gen~2) дало прирост $+11.8\%$ для MLP. Введение survivability heatmaps (Gen~3) обеспечило lookahead внутри раунда без явного дерева перебора. Каждое поколение превзошло предыдущее при неизменном алгоритме PPO.

  \item \textbf{Порочный круг «слабый критик --- шумные признаки» требует curriculum.} CNN не способна обучиться с нуля на тензоре $(14,8,8)$: топологически информативные каналы (survivability, blob, dead zones) бесполезны при случайной политике. Transfer learning из Gen~2 разорвал порочный круг и позволил воспользоваться новыми каналами, фактически реализовав неявный curriculum learning.

  \item \textbf{Behavior Cloning решает проблему холодного старта.} На поле $16\times16$ прямое PPO-обучение не даёт полезного сигнала: случайный агент набирает лишь ${\approx}0.5$ lines/эп. Предобучение на $500\,000$ демонстрациях эвристики выводит политику из «мёртвой зоны»: агент стартует с ${\approx}4$--$5$ lines/эп. (против ${\approx}0.5$ у случайного), а PPO fine-tuning улучшает его до $18.76$ lines --- значительно выше уровня учителя ($7.21$).

  \item \textbf{Нейросетевые агенты превзошли эвристику по всем ключевым метрикам.} На $8\times8$ превосходство умеренное ($+4$--$32\%$ по lines). На $16\times16$ разрыв составил $2.6\times$ ($18.76$ vs $7.21$ lines), Cohen's $d=1.39$ --- очень большой эффект. На пороге $\geq30$ lines нейросети превосходят эвристику в $13\times$. Нейросеть самостоятельно выработала диагональную стратегию заполнения --- нелокальную политику, недостижимую одношаговой жадной эвристикой.

  \item \textbf{CNN и ViT статистически неразличимы при равном бюджете параметров (${\approx}1.25$М).} На $16\times16$ оба агента показали ${\approx}18.5$ lines ($p=0.42$, Mann--Whitney~U). Ограничивающим фактором является качество BC+PPO пайплайна, а не выбор архитектуры.
\end{enumerate}

\paragraph{Ограничения.}

\begin{itemize}
  \item \textbf{PPO не сходится на $16\times16$ без BC.} Прямое RL-обучение на большом поле застревает ниже $1$ line/эп. неопределённо долго --- метод зависит от качества эксперта: слабая эвристика дала бы слабый warm-start.
  \item \textbf{ViT не принёс преимущества над CNN.} Несмотря на теоретически более подходящее индуктивное смещение, результаты CNN и ViT статистически неразличимы. Возможно, $16$ патчей недостаточно для выявления нелокальных паттернов, или $15$М шагов PPO мало для реализации потенциала трансформера.
  \item \textbf{Оценка проводилась на фиксированном \texttt{seed=42}.} $500$ эпизодов обеспечивают достаточную статистическую мощность для обнаруженных размеров эффектов (Cohen's $d=1.39$), однако перекрёстная валидация по нескольким seeds остаётся направлением для дополнительной проверки воспроизводимости.
  \item \textbf{Эвристика геометрически безнадёжных ситуаций не преодолевается.} При раскладе, где все три фигуры несовместимы с оставшимся пространством, ни одна стратегия не спасает --- нижняя граница числа lines определяется случайностью пула.
\end{itemize}

% ════════════════════════════════════════════════════════════════
%  Список литературы
% ════════════════════════════════════════════════════════════════
\newpage
\begin{thebibliography}{99}

\bibitem{sutton2018}
  Sutton R.S., Barto A.G.
  \textit{Reinforcement Learning: An Introduction}.
  2nd ed. MIT Press, 2018.

\bibitem{schulman2017}
  Schulman J., Wolski F., Dhariwal P., Radford A., Klimov O.
  Proximal Policy Optimization Algorithms.
  \textit{arXiv:1707.06347}, 2017.

\bibitem{schulman2015}
  Schulman J., Moritz P., Levine S., Jordan M., Abbeel P.
  High-Dimensional Continuous Control Using Generalized Advantage Estimation.
  \textit{arXiv:1506.02438}, 2015.

\bibitem{mnih2015}
  Mnih V. et al.
  Human-level control through deep reinforcement learning.
  \textit{Nature}, 518(7540):529--533, 2015.

\bibitem{huang2020}
  Huang S., Ontañón S.
  A Closer Look at Invalid Action Masking in Policy Gradient Algorithms.
  \textit{arXiv:2006.14171}, 2020.

\bibitem{gymnasium2022}
  Towers M. et al.
  Gymnasium.
  \textit{Zenodo}, 2023.
  \url{https://github.com/Farama-Foundation/Gymnasium}

\bibitem{sb3}
  Raffin A. et al.
  Stable-Baselines3: Reliable Reinforcement Learning Implementations.
  \textit{Journal of Machine Learning Research}, 22(268):1--8, 2021.

\bibitem{ng1999}
  Ng A.Y., Harada D., Russell S.
  Policy Invariance Under Reward Transformations: Theory and Application to Reward Shaping.
  \textit{ICML}, 1999.

\bibitem{dellacherie2003}
  Dellacherie P.
  Designing the perfect Tetris AI.
  \textit{Online resource}, 2003.
  \url{https://codemyroad.wordpress.com/2013/04/14/tetris-ai-the-near-perfect-player/}

\bibitem{pomerleau1991}
  Pomerleau D.A.
  Efficient Training of Artificial Neural Networks for Autonomous Navigation.
  \textit{Neural Computation}, 3(1):88--97, 1991.

\bibitem{ross2011}
  Ross S., Gordon G., Bagnell D.
  A Reduction of Imitation Learning and Structured Prediction to No-Regret Online Learning.
  \textit{AISTATS}, pp.~627--635, 2011.

\bibitem{dosovitskiy2020}
  Dosovitskiy A. et al.
  An Image is Worth $16\times16$ Words: Transformers for Image Recognition at Scale.
  \textit{ICLR}, 2021.

\bibitem{hafner2023}
  Hafner D. et al.
  Mastering Diverse Domains through World Models.
  \textit{arXiv:2301.04104}, 2023.

\end{thebibliography}

\end{document}
